\providecommand{\devnum}[1]{#1}

\section{Scheduling Algorithm}
\label{sec:scheduling}

We build the scheduler in two layers, because the ranking that shortens the queue is also what
can leave one submission behind without end. A base policy orders the queue by a score
computed at arrival, and a wrapper limits how much work may overtake any waiting job.
Because the two are independent, the wrapper accepts any base policy, and its guarantee does
not depend on that policy behaving well.

\subsection{Ranking by Expected Cost}
\label{sec:rank}

When a server frees, the base policy takes the waiting job with the smallest score.
For a batch on one server, the pairwise-interchange argument orders jobs by their
conditional mean service times. However, the same conclusion does not hold in general on $k\geq2$ servers. For example, on two servers let $A=2$, $C=3$, and let $B$ be
$0.1$ with probability $0.9$ and $100$ with probability $0.1$. Ordering the three
conditional means starts $A$ and $C$ and has expected mean wait $2/3$, while starting $A$ and $B$ instead has expected mean wait $0.096667$. Hence expected-cost ordering is an empirically motivated base policy here, not a multiserver optimality
claim.
% Source for the two-server calculation and its independent script check:
% docs/referee_readthrough/independent_review/checks_output.txt, CHECK 1
% (script: docs/referee_readthrough/independent_review/checks.ps1).

Both raw-scale squared error and the Tweedie mean deviance have the conditional mean
$\mathbb{E}[C_i\mid x_i]$ as their unrestricted population optimum. Their fitted orders can nevertheless differ with finite model capacity, the link function,
regularisation, sample weighting and optimisation. We therefore compare objectives
empirically while keeping the features, model size and rolling-origin protocol of
Section~\ref{sec:prediction} fixed. For SPJF-E we use raw seconds with a Tweedie loss of power $1.5$, which accommodates mass near zero and a long right tail
\cite{jorgensen1987edm,yang2018tweedieboost}. SPJF-log uses the same pipeline fitted
by squared error on $\log(1+C)$, and ranks by that fitted value, which orders jobs as
its inverse would.
% Source for the fitted Tweedie power and objective comparison:
% evidence/codebench_service_v2/out_service_v2.txt, SUMMARY 1.

The literal inverse $\exp(\mathbb{E}[\log(1+C)\mid x])-1$ is, by Jensen's inequality, no
larger than the conditional arithmetic mean, with equality only in the degenerate case. The gap is not determined by the conditional variance alone, and to use a log-scale
regression as a
mean estimator under its stated conditions, one needs a residual-distribution
correction such as the smearing estimator of Duan~\cite{duan1983smearing}.

We use only the order of the fitted score. Its level does not transfer across periods, since the share of heavy jobs in the training periods differs from the target
period by a factor of several, so we never read the score as a quantity, neither as an
estimated finish time nor as a threshold. Hence the scheduler is invariant to a monotone recalibration of the score.

Section~\ref{sec:exp_rank} measures what the fitted objective gains, together with
the alternatives we tried and rejected: a predicted quantile, a two-class rule, and
squared error on raw seconds. Its partial-oracle repair replaces selected scores with truth and is not invariant to an arbitrary monotone rescoring of the original output.
Thus it identifies the errors of this fitted predictor and does not establish a
universal asymmetry of prediction errors or a different population estimand.

No choice of ranking function bounds the worst case. Whatever the score is and however it was fitted, a job the score gets wrong can still be passed without end, and only the guard of Section~\ref{sec:guard} bounds it.

\subsection{The Overtake-Budget Guard}
\label{sec:guard}

We fix a base policy $A$, and for every waiting job $q$ and every time $t$ we define
\[
  \mathrm{over}[q](t) \;=\;
  \sum \{\, C_j \;:\; j > q \text{ in rank, } j \text{ has completed by } t \,\},
\]
the total true service of the higher-ranked jobs that have finished by $t$. Such a job
arrived at or after $a_q$, so it necessarily ran while $q$ waited. We allow job $q$ a budget $\mathrm{budget}(q,t)$ on that quantity, and the fired set at a
dispatch epoch is
$E(t) = \{\, q \text{ waiting} : \mathrm{over}[q](t) \ge \mathrm{budget}(q,t) \,\}$.
At every dispatch epoch, that is, at every instant at which a server is free and the
queue is non-empty, the wrapper serves the minimum-rank member of $E(t)$ when $E(t)$
is non-empty and otherwise serves whatever $A$ chooses (Figure~\ref{fig:guard}).

Serving the minimum-rank member matters when budgets differ between jobs, because it guarantees that whenever a job $j$ is dispatched ahead of a lower-ranked waiting job
$i$, the job $i$ was not in $E(t)$, hence $\mathrm{over}[i](t) < \mathrm{budget}(i,t)$
at that very instant. The per-job bound of Section~\ref{sec:theory} rests on that
step. With a constant budget $\mathrm{budget}(q,t) \equiv B$ the quantity
$\mathrm{over}[q](t)$ is non-increasing in rank, so $E(t)$ is non-empty exactly when
the queue head belongs to it and the rule degenerates to a test on the head, and setting $B = 0$ recovers FCFS.

\paragraph{The cap carries the promise; the shape does not.}
As Section~\ref{sec:theory} shows, the guarantee reads the budget only through an
upper bound on it. If $0 \le \mathrm{budget}(q,t) \le B_{\max}$ for every waiting job and
every instant, then every job waits at most
$W_{\mathrm{FCFS}} + B_{\max}/k + (3-2/k)L$, for every base policy, every arrival
sequence and every set of predictions. The budget below the cap is a design freedom, which decides how much of the base
policy's benefit survives and which jobs bear the harm and does not change the promise. We study the one-line family
\[
  \mathrm{budget}(q,t) \;=\;
  \min\bigl(\,B_0 + \gamma\,n_q + \eta\,k\,(t - a_q),\ B_{\max}\,\bigr),
\]
with $B_0, \gamma \ge 0$, $0 \le \eta < 1$ and $n_q$ the number of jobs waiting when
$q$ arrives, and we read three shapes off it, \emph{constant} ($\gamma = \eta = 0$),
\emph{age-relative} ($\eta > 0$), and \emph{queue-length} ($\gamma > 0$). The
coefficient $\eta$ says how much of the system's capacity, over a job's own waiting
time, may go to jobs that arrived after it. In other words, the allowance grows
with the wait a job has already had. We chose a fraction of capacity rather than a
number of seconds because it is dimensionless, so it needs no retuning when $k$ or the
load changes. The term $\gamma n_q$, by contrast, widens the allowance only for jobs that
arrive into a queue that is already long, and it is flat in the wait. All three of $B_0$, $\gamma$ and $B_{\max}$ are measured in work and scale with $k$.

Only the age-relative shape carries a second, multiplicative bound, and only when
$\gamma = 0$: $W < (W_{\mathrm{FCFS}} + B_0/k + (3-2/k)L)/(1-\eta)$. With
$\gamma > 0$ the constant part of the budget differs from job to job. Hence the multiplicative form holds per job, with that job's own constant, and the operator
has no single number to state in advance (Supplementary Section~S4.6), so for $\gamma > 0$ we claim the additive bound alone.

\begin{algorithm}
\caption{Overtake-budget guard wrapped around a base policy $A$}
\label{alg:guard}
\begin{algorithmic}[1]
\State \textbf{parameters:} cap $B_{\max} \ge 0$; shape $B_0, \gamma \ge 0$, $\eta \in [0,1)$; base policy $A$
\State \textbf{state:} waiting set $Q \gets \emptyset$; $\mathrm{over}[q] \gets 0$ for every job $q$
\While{some job has not started}
  \State $t \gets$ next event time
  \ForAll{jobs $c$ completing at $t$} \Comment{$C_c$ becomes known here and nowhere earlier}
    \ForAll{$q \in Q$ with $q < c$ in rank}
      \State $\mathrm{over}[q] \gets \mathrm{over}[q] + C_c$
    \EndFor
  \EndFor
  \ForAll{jobs $i$ arriving at $t$}
    \State $n_i \gets |Q|$; \ $Q \gets Q \cup \{i\}$; \ $\mathrm{over}[i] \gets 0$ \Comment{$n_i$ fixes $i$'s constant part}
  \EndFor
  \While{a server is free \textbf{and} $Q \neq \emptyset$} \Comment{dispatch epoch}
    \State $E \gets \{\, q \in Q : \mathrm{over}[q] \ge \min(B_0 + \gamma n_q + \eta k (t - a_q),\, B_{\max}) \,\}$
    \If{$E \neq \emptyset$}
      \State $j \gets$ the minimum-rank member of $E$ \Comment{guard fires}
    \Else
      \State $j \gets A(Q, t)$ \Comment{base policy decides}
    \EndIf
    \State dispatch $j$ on the free server; \ $W[j] \gets t - a_j$; \ $Q \gets Q \setminus \{j\}$
  \EndWhile
\EndWhile
\end{algorithmic}
\end{algorithm}

% Figure: the guard as a wrapper, and the accounting it runs.  Everything drawn
% here is Algorithm~\ref{alg:guard} and the definitions above it; the lower panel
% is a schematic of over[q](t) against budget(q,t), not a measurement.
\begin{figure}[htbp]
\centering
\resizebox{\textwidth}{!}{%
\begin{tikzpicture}[x=1cm,y=1cm,font=\footnotesize,
   box/.style={draw, rounded corners=1pt, align=center, inner sep=3pt},
   dec/.style={draw, align=center, inner sep=3pt},
   ar/.style={-{Latex[length=4pt]}}]

  \node[box]           (queue) at (0,0)    {waiting queue $Q$\\[-1pt]\scriptsize ranks, scores $\hat{C}$};
  \node[dec, anchor=west] (test)  at (3.0,0)  {$E(t)=\emptyset$?\\[-1pt]
        \scriptsize $E(t)=\{q:\mathrm{over}[q]\ge\mathrm{budget}(q,t)\}$};
  \node[box, anchor=west] (base)  at (9.3,1.25) {base policy $A$\\[-1pt]\scriptsize any ranking, learned or not};
  \node[box, anchor=west] (fire)  at (9.3,-1.25){guard fires\\[-1pt]\scriptsize min-rank member of $E(t)$};
  \node[box, anchor=west] (disp)  at (13.9,0) {dispatch on\\the free server};
  \node[box, anchor=east] (comp)  at (2.4,-2.6) {completion of $j$\\[-1pt]\scriptsize $C_j$ learned here, not earlier};

  \draw[ar] (queue.east) -- (test.west);
  \draw[ar] (test.east) -- ++(0.35,0) |- node[pos=0.72,above,font=\scriptsize] {yes} (base.west);
  \draw[ar] (test.east) -- ++(0.35,0) |- node[pos=0.72,below,font=\scriptsize] {no}  (fire.west);
  \draw[ar] (base.east) -- ++(0.35,0) |- (disp.west);
  \draw[ar] (fire.east) -- ++(0.35,0) |- (disp.west);
  \draw[ar] (comp.east) -- (test.south |- comp.east) -- (test.south);
  \node[anchor=west,font=\scriptsize] at ([xshift=3mm]test.south |- comp.east)
        {$\mathrm{over}[q] \mathrel{+}= C_j$ for every waiting $q$ of lower rank};
\end{tikzpicture}}

% Referee item G11 asked for Figures 1 and 2 to be set side by side.  They sit in
% different sections eleven pages apart, so that is not available; the vertical extent of
% this figure is tightened instead, here and in the lower panel's width.
\vspace{1mm}

\resizebox{0.46\textwidth}{!}{%
\begin{tikzpicture}[x=1cm,y=1cm,font=\footnotesize,
   ar/.style={-{Latex[length=4pt]}}]
  \draw[ar] (0,0) -- (6.4,0) node[right] {$t$};
  \draw[ar] (0,0) -- (0,3.4);
  \draw[densely dashed] (0,1.1) -- (3.1,2.7) -- (6.0,2.7);
  \node[anchor=west] at (3.3,3.15) {$\mathrm{budget}(q,t)$};
  \node[anchor=east] at (-0.05,1.1) {$B_0$};
  \node[anchor=east] at (-0.05,2.7) {$B_{\max}$};
  \draw[thick] (0,0) -- (1.0,0) -- (1.0,0.7) -- (2.2,0.7) -- (2.2,1.5)
               -- (3.4,1.5) -- (3.4,2.2) -- (4.6,2.2) -- (4.6,2.8) -- (5.2,2.8);
  \node[anchor=west] at (0.15,2.35) {$\mathrm{over}[q](t)$};
  \draw[ar] (5.6,1.75) -- (4.68,2.72);
  \node[anchor=west] at (5.5,1.55) {guard fires};
\end{tikzpicture}}
\caption{The guard of Algorithm~\ref{alg:guard} wrapped around a base policy
(above) and the accounting it runs for one waiting job (below). Only completed
work is charged, because $C_j$ is unknown until $j$ finishes, so
$\mathrm{over}[q]$ is a step function with a step at each completion of a
higher-ranked job; the first dispatch epoch at which the step function reaches
the budget sends $q$ next. The budget drawn is the age-relative shape, rising at
rate $\eta k$ from $B_0$ and capped at $B_{\max}$; the guarantee depends on the
cap alone, so any curve under the dashed ceiling carries the same promise. The
lower panel is schematic. Sources: the definitions in Section~\ref{sec:guard};
$B_{\max} = k(G - (3-2/k)L)$ is Corollary~\ref{cor:params}.}
\label{fig:guard}
\end{figure}


\paragraph{Parameters from a service promise.}
The operator states $G$, the number of seconds by which any job may exceed its FCFS
waiting time, and reads off $B_{\max} = k\bigl(G - (3 - 2/k)\,L\bigr)$, the inversion
of the uniform certificate proved in Section~\ref{sec:theory}, and no data are needed for this step. The rule issues no positive budget below $(3-2/k)L$, while the zero budget remains
admissible and reproduces FCFS exactly. In Section~\ref{sec:theory} we also give an $L$-scale floor, but only for this wrapper family when it must protect
against every opaque base-policy proposal sequence. The promise $G$ follows from $B_{\max}$ alone, while the shape parameters $B_0$, $\gamma$ and $\eta$ decide how much of
the base policy's benefit survives and which jobs pay for it, so we choose them on a
validation trace that excludes the test period. That pool is held out of the selection
and not of the development process, and Section~\ref{sec:exp_prov} states what the
selection therefore does and does not establish. We fix the rule before the grids are run, state it in full in Supplementary
Section~S5, and give the grids and what they selected in Supplementary Table~S4.


\paragraph{What is charged, and why no service time is needed at dispatch.}
We charge the budget with the true executed work of overtakers, and only when they
complete. At the instant the guard decides, it has read $C_j$ only for jobs that have
already finished, and a job still in service contributes nothing. That is what makes our wrapper implementable in a system that learns a job's cost at its completion, and
it is also why the bound carries an additive term in place of an exact
accounting. When the last overtaker of $i$ is dispatched, up to $k-1$ further
overtakers may be in service with their work not yet charged, and the overtaker being
dispatched is not charged either, so the work that actually passes $i$ can exceed the
budget by at most $kL$. Under the timeout at $L$, executed work means $\min(C_j, L)$, and because the log records that quantity and the guard
charges it, no assumption on raw
sizes enters.


\paragraph{Cost per event.}
We write Algorithm~\ref{alg:guard} for readability rather than for execution, since
its inner loop over waiting jobs at each completion is linear. With a Fenwick tree over completed work and a segment tree over the window of waiting
ranks, we make each event cost $O(\log n)$. We carry all budget arithmetic in integer microseconds, so the firing
test is exact and no floating-point clock enters the schedule. We give the two structures and the rewriting of the firing test in Supplementary
Section~S5. Thus a promise in seconds becomes a number the dispatcher can check, and the dispatcher never reads
a prediction to check it.

\paragraph{Why the budget is measured in work and not in positions.}
Earlier bounded-overtaking mechanisms count passes~\cite{grosof2021nudge,vanhoudt2022nudgek}
or limit service to the oldest jobs~\cite{grosof2022wcfs} instead. A count is loose
against a work budget by the ratio $L/\bar{C}$, \devnum{150} on the main trace ($L = 60$ s
against a mean cost of \devnum{0.40} s). We give the argument in Supplementary
Sections~S4.4 and~S4.9, and
Section~\ref{sec:experiments} measures what the looseness costs on a real load.
% Source: prechecks/cost_ratio/out_mean_cost_ratio.txt (L / mean cost = 150.1, mean 0.3997 s).

\subsection{Compared Policies}

In Table~\ref{tab:policies} we name the policies evaluated in
Section~\ref{sec:experiments} and fix the names used in every table of this paper, and Supplementary Section~S5.7 states what each one is for at length.
Our method is Guard($G$) wrapped around SPJF-E, and the others separate its parts and
test it.

\begin{table}[htbp]
\centering
\caption{Policies compared in the experiments, and the name each carries throughout.}
\label{tab:policies}
\small
\begin{tabular}{p{0.19\textwidth}p{0.71\textwidth}}
\hline
Name & Role \\
\hline
FCFS & Deployed baseline, and the reference the guarantee is stated against \\
SJF & Undeployable reference that sorts by true service time; the denominator of the gap closed \\
SPJF-E & Ranks by the expected-cost score of Section~\ref{sec:rank}, unguarded \\
SPJF-log & The same pipeline on the log scale, and the other candidate scores of Section~\ref{sec:rank} \\
Guard($G$) & SPJF-E wrapped in the overtake-budget guard at the promise $G$. This is the method \\
Guard-fixed($G$), Guard-age($G$), Guard-queue($G$) & The same wrapper with the shape the rule picks inside one grid. Ablation rows \\
Fixed($G$) & The same wrapper at the equal-promise budget $\mathrm{budget} \equiv B_{\max}$; infeasible under the harm rule \\
Skip($G$) & Position-count guard, finite-skip family, counted at dispatch (Supplementary Section~S4.4) \\
Corrupted scores & SPJF-E and Guard($G$) under a reversed order, a random order and an inverted-tail score \\
Counterexamples & Retained traces from two audits and adversarial review, checked against the bounds \\
\hline
\end{tabular}
\end{table}

In every simulation that includes the guard we assert the per-job bound of
Section~\ref{sec:theory} job by job inside the run
(Section~\ref{sec:exp_bound}).
